\section{AI 编程与 Agent 工具链：从“补全代码”到“管理规范”}

\subsection{工作模式的变化：从写每一行到定义每一步}

访谈中关于 coding agents 的讨论很具体。开发者体验从“自动补全”转向“多轮计划 + 局部执行 + 人类审查”。这使高级工程师受益更大，因为他们更擅长写清楚规格、识别隐含约束、快速定位失败模式。

\begin{importantbox}{为何资深开发者更容易把 AI 用好}
\begin{itemize}[leftmargin=1.5em]
    \item 能提前给出边界条件和非功能需求（性能、稳定性、安全）；
    \item 知道哪些模块可并行、哪些必须串行；
    \item 对错误模式有先验，能更快给出修复方向而非重复尝试。
\end{itemize}
\end{importantbox}

\subsection{“Spec-driven development”成为核心技能}

Nathan 提到很多失败并不是模型不会写代码，而是规格不充分。工具会机械执行用户意图的字面形式，不会自动补足业务语义。于是“写出可执行规格”从产品经理技能外溢到每个工程角色。

\begin{knowledgebox}{面向 Agent 的规格模板（可落地）}
一个高质量任务描述通常包含：目标输出、边界条件、禁止事项、可用工具、验收标准、失败回滚方案。这个模板本质上是把 tacit knowledge 显式化，减少“模型猜你意思”的空间。
\end{knowledgebox}

\begin{warningbox}{全自动编程的主要瓶颈仍是系统复杂性}
访谈中有一句很实在的话：模型会把同一个错误命令重复执行很多次。它暴露的不是“不会写代码”，而是“在复杂系统里缺少稳健的问题求解策略”。在分布式系统、遗留系统、跨服务变更里，这个问题尤其明显。
\end{warningbox}

\subsection{本章小结}

AI 编程在提效上已是确定趋势，但“谁定义规范、谁做最终审查”依旧是决定工程质量的核心。Agent 时代对人的要求不是更少，而是更高阶。

